\section{Evaluation Protocol and Implementation Status}
\label{sec:evaluation}

The complete benchmark matrix is intentionally left for the final clean rerun.
This section fixes the questions, corpus, metrics, budgets, and validation chain
implemented in the artifact; it reports only results already reproduced from a
pinned clean revision.

\subsection{Research Questions}

\begin{description}
\item[RQ1: Atomic compilation.] Does MOOSE evidence compile into a closed,
lifted AgentSpeak library whose branches satisfy the declared certificates?
\item[RQ2: Symbol invariance.] Do compiler decisions remain unchanged under
predicate/action renaming, compatible parameter permutation, and irrelevant
fluent injection?
\item[RQ3: Temporal composition.] Does each supported DFA guard compile into a
preservation-certified controller without changing order across DFA edges?
\item[RQ4: Executable validity.] Does Jason produce a complete primitive-action
trace that VAL accepts for the corresponding PDDL instance?
\end{description}

\subsection{Benchmark Scope}

The corpus contains all MOOSE companion domains with direct nonempty train/test
splits, plus project-added serialized-width families. The groups in
Table~\ref{tab:benchmarks} describe evaluation properties; they do not route a
domain to a different compiler.

\begin{table}[t]
\centering
\caption{Implemented benchmark corpus.}
\label{tab:benchmarks}
\begin{tabular}{p{0.26\columnwidth}p{0.63\columnwidth}}
\toprule
Property group & Domains \\
\midrule
Classical & Barman, Ferry, Gripper, Logistics, Miconic, Rovers, Satellite,
Transport \\
Numeric & Numeric Ferry, Numeric Miconic, Numeric Minecraft, Numeric Transport \\
Serialized width & Blocksworld Clear, Blocksworld On, Blocksworld Tower, Depots \\
\bottomrule
\end{tabular}
\end{table}

MOOSE domains use the official companion \texttt{training/} and
\texttt{testing/} directories. Blocksworld Clear and On use the published
no-constants QClear/QOn splits. Blocksworld Tower uses the first quarter of its
ordered instances for training; Depots uses the first half because its source
contains only 22 instances. These deterministic project splits are fixed before
execution and recorded in source metadata.

\subsection{Budgets and Systems}

MOOSE-native settings follow its paper where applicable: three goal
permutations by default, five synthesis repetitions, a 12-hour/32-GB synthesis
budget, and a 30-minute/8-GB test-planning budget
\cite{Chen2026Moose}. The artifact keeps the 12-hour synthesis timeout but uses
a 16-GiB external-process safety cap; this memory difference is a declared
reproduction deviation. MOOSE, Jason, and validation worker counts are system
parameters rather than generalized-planning hyperparameters. The batch scripts
default to six workers, a 30-minute Jason timeout, a 30-minute VAL timeout, and
a 64-MiB Java thread stack.

The implemented toolchain uses the MOOSE readable-policy exporter, Clingo for
branch selection, \texttt{ltlf2dfa} with a real MONA binary, Jason for
AgentSpeak execution, and Dockerized VAL for independent PDDL validation. Every
run records its Git commit and separately records tracked modifications and
untracked files. Results from a tracked-dirty tree are diagnostic and are not
eligible for the final matrix.

\subsection{Metrics and Validation}

For each domain we record evidence parsing, atomic compilation, selected branch
count, context/body size, and closure status. For each test problem we record
Jason success, timeout or failure, primitive action count, VAL success, and
wall-clock time. A problem counts as solved only when Jason emits its explicit
success marker and VAL accepts the complete exported trace. Prefix actions from
a failed or timed-out intention are diagnostic and are not counted as a plan.

For temporal inputs we additionally record DFA construction, supported-fragment
validation, transition count, serialization certificate, negative-guard
preservation, and final finite-trace acceptance. Classical PDDL conjunctions
are used as one-transition achievement guards for infrastructure validation;
they are not presented as natural-language \ltlf predictions.

The symbol-invariance suite uses synthetic domains whose vocabulary is renamed
and whose irrelevant predicates are changed. The staged-preparation regression,
for example, uses fabrication and tool predicates rather than any Rovers
vocabulary while preserving the same PDDL dependency structure. This tests the
claim that compiler rules depend on schemas rather than benchmark names.

\subsection{Current Verified Checkpoint}

At commit \texttt{ac934327}, the full Python suite contains 340 passing tests;
the only warnings are two deprecations in the third-party Lark package. Ruff
passes for the complete repository. A clean-revision Rovers run compiles one
atomic library and obtains Jason- and VAL-valid traces for all 90 official test
instances, including the two staged image-production regressions that motivated
existential preparation projection. A clean Blocksworld-On regression produces
a 30-action VAL-valid trace, confirming that the functional-fluent restriction
does not retain the earlier over-broad branch reordering.

These checks establish implementation behavior for the reported cases, not the
final cross-domain coverage claim. The final paper table will be generated only
after rerunning all domains from one clean pinned revision with the budgets
above.

\subsection{Reproducibility Artifacts}

The repository contains a benchmark registry, pinned source metadata, canonical
domain libraries, per-test AgentSpeak programs, exported PDDL traces, VAL
records, DFA diagnostics, and timestamped summary JSON. The final-paper script
regenerates result macros from these artifacts rather than transcribing numbers
by hand. The submission package will include the official AAAI reproducibility
checklist after the references.
